\lhead{\textbf{Basic Algorithms, Fall 2025 \\ CSCI-UA.0310-001}}
\chead{\Large{\textbf{Homework 5}}}
%%%%%%%%%%%%%%%%%%%%%%%%%%%%%%%%%%%%%%%
% ENTER NAME BELOW!
%%%%%%%%%%%%%%%%%%%%%%%%%%%%%%%%%%%%%%%
\rhead{\textbf{Professor Regev}\\\textbf{Name: } ~~~~~~~~~~~~~~~~~~~~~~~~~~~~~~~~~~~~~}
%REPLACE THE TILDES WITH YOUR NAME
\runningheadrule
\firstpageheadrule
\cfoot{}

\section*{Due October 13 (11:59 p.m.)}
\intro

\subsection*{1-0 List all your collaborators and sources: (\texorpdfstring{$-\infty$}{-∞} points if left blank)}
\vspace{0.75cm}



\section*{Problem 5-1 -- Counting Sort (20 points)}

Recall that in Counting (or Small-Key Sort) sort, we append the current element $A[i]$ to the end of the list $L_{k_i}$ where $k_i \in \{1,\dots,k\}$ is the key associated with $A[i]$. Consider instead adding $A[i]$ to the beginning of the list (the rest of the algorithm is unchanged and outputs the concatenated final lists $L_1,\dots,L_k$). What property of Counting sort is broken? 

\begin{solution}

\vspace{0.5cm}

The modified procedure \emph{breaks stability}.

% \paragraph{Definition (Stability).} A sorting algorithm is \emph{stable} if for any indices $i<j$ with $k_i = k_j$ (i.e., equal keys), the element $A[i]$ appears before $A[j]$ in the output.

In the standard (stable) implementation of Counting/Small-Key sort, 

each element $A[i]$ is \emph{appended to the end} of its bucket $L_{k_i}$

\vspace{0.2cm}


Therefore, within every bucket $L_{k}$, elements keep their original left-to-right order

Concatenating $L_1, L_2, \dots, L_k$ thus preserves the relative order of items with equal keys

\vspace{0.5cm}

If instead each $A[i]$ is \emph{inserted at the beginning} of $L_{k_i}$, 


\hspace{1cm}
then within every bucket the arrival order is \emph{reversed}

Consequently, when the buckets are concatenated, the elements sharing the same key appear in the \emph{reverse} of their input order, violating stability

\vspace{0.5cm}

\textbf{Counterexample}


\vspace{0.5cm}


Let the input be $A=[(a,2),(b,1),(c,2)]$, where the second component is the key


\quad Standard (append-to-end) buckets: 

\hspace{1.5cm}
$L_1=[b],; L_2=[a,c]$, output $[(b,1),(a,2),(c,2)]$ , which is stable


\quad Modified (prepend-to-front) buckets: 

\hspace{1.5cm}
$L_1=[b],; L_2=[c,a]$, output $[(b,1),(c,2),(a,2)]$, 

\hspace{1.5cm}
where $a$ and $c$ (same key) appear in reversed order

\vspace{0.5cm}

Hence, stability is lost

\end{solution}

% \newpage

\section*{Problem 5-2 -- Sorting Efficiently (30 Points)}

\begin{enumerate}
	\item (15 points)
	Suppose you have an array $A$ of length $2^{16}$, where each entry is a 128-bit integer. For each of \emph{Radix Sort}, \emph{Counting Sort}, \emph{Merge Sort}, state the optimal runtime for each algorithm running on $A$ (Note that for radix sort, this means choosing the optimal base -- please state this concretely.) Write the runtimes as an approximate number of steps — you may assume Merge Sort runs in time $n\log n$ (i.e., no constants), and Counting Sort runs in $2n+2k$ steps. Finally, state which of these algorithms is most efficient for sorting $A$.
	
	\begin{solution}
		\vspace{0.5 cm}



        \begin{enumerate}
        
  % \item \textbf{Parameters.} Let \(n=2^{16}=65536\). Each key is a 128-bit integer (key-space size upper bound \(k=2^{128}\)).

  \item \textbf{Merge Sort}
  \[
    T_{\text{merge}}(n)=n\log n = 2^{16}\cdot 16 = 65536\times 16 = 1048576 \text{ steps}
  \]

  \item \textbf{Counting Sort}
  \[
    T_{\text{count}}(n,k)=2n+2k
  \]
  For 128-bit keys taken at full range \(k=2^{128}\),
  
  \[
    T_{\text{count}} \approx 2\cdot 2^{128} = 2^{129} \text{ steps}
  \]

  \item \textbf{Radix Sort} 
  
  Use base \(b\) per pass (so each pass uses Counting Sort on \(b\) buckets)

  
  Number of passes \(d=\dfrac{128}{\log_2 b}\)
  
  Each pass costs \(2n+2b\)
  
  Thus
  \[
    T_{\text{radix}}(b)=\frac{128}{\log_2 b}\cdot(2n+2b)
  \]
  A natural (near-optimal) choice is \(b=2^{16}=65536=n\)
  
  Then
  \[
    \log_2 b=16,\quad d=\frac{128}{16}=8,
  \]
  per-pass cost \(2n+2b=2\cdot65536+2\cdot65536=262144\), and
  \[
    T_{\text{radix}}(b=2^{16})=8\times 262144 = 2097152 \text{ steps}
  \]
\end{enumerate}

\textbf{Comparison and conclusion}
  \[
    T_{\text{merge}}=1048576,\quad
    T_{\text{radix (best)}}\approx 2097152,\quad
    T_{\text{count (full 128-bit range)}}\approx 2^{129}\ 
  \]
  Therefore \textbf{Merge Sort} is the most efficient for sorting this array under the given assumptions

        
	\end{solution}
    
    \item (15 points)
	Assume you are given an array of $n$ integers with many
	duplications, so that you know that there are at most $\log n$
	distinct elements in the array. Show how to sort this array in time
	$O(n \log\!\log n)$.
	
\begin{solution}
\vspace{0.5 cm}

% We are given an array of $n$ integers containing at most $\log n$ distinct values. We can sort it in $O(n \log\!\log n)$ time as follows:

\begin{enumerate}
    \item Traverse the array and use a hash table or dictionary to count the frequency of each distinct element
    
    This takes $O(n)$ time

    
    \item Extract the distinct elements (at most $\log n$ of them) and sort them using any comparison-based sorting algorithm
    
    Sorting $k = O(\log n)$ elements takes $O(k \log k) = O(\log n \cdot \log\!\log n)$ time
    
    \item Finally, iterate over the sorted distinct elements, 
    
    and for each element, write it back into the output array according to its frequency
    
    This reconstruction step takes $O(n)$ time

    
\end{enumerate}

The total running time is
\[
O(n) + O(\log n \cdot \log\!\log n) + O(n) = O(n) 
\]

since big $\mathcal{O}$ notation is the upper bound

The array can be sorted in $O(n \log\!\log n)$ time
\end{solution}





\end{enumerate}

% \newpage

\section*{Problem 5-3 -- Pumpkin Pickin' (50 points)}

% Consider a neighborhood of $n$ houses in a line, represented by an array $A$,
% where $A[i]$ is the value of the items that can be stolen from the $i^{\text{th}}$ house.
% A dastardly robber would like to steal as much as possible from this neighborhood; however, if they attempt to rob two \emph{adjacent} houses, this will trigger the neighborhood's collective alarm system and call the police. In this problem, you will use dynamic programming to give an efficient algorithm
% determining the maximum total value of items the robber can steal without triggering the alarm.

You work at a pumpkin patch and you are tasked with harvesting the most pounds of pumpkin without ruining the aesthetic of the pumpkin vine.  In particular, you are not to harvest two \emph{adjacent} pumpkins.  The vine of $n$ pumpkins is represented by an array $A$ where $A[i]$ is the weight of pumpkin $i$ (in pounds). Use dynamic programming to give an efficient algorithm determining the maximum total weight of pumpkin you can harvest while adhering to the rule of not harvesting two in a row. 

\begin{enumerate}
    \item (8 points)
    Define the subproblems
    for the dynamic programming algorithm.
    How many subproblems are there?
    \begin{solution}
         
        \vspace{0.4cm}

        The subproblem is 

        \[
        \text{the maximum weight that we can achieve with the first } i \text{ pumpkins, i.e.}A[1:i]
        \]
        
        The number of sub-problems is $n$



        
    \end{solution}
    \item (8 points) Write out the base cases (there should be two of them!).
    \begin{solution}
        \vspace{0.2cm}

        if $i = 1$:
        
        \hspace{1cm} return $A[1]$
        
        elif $i = 2$:
        
        \hspace{1cm} return max$(A[1],A[2])$


        
    \end{solution}
    \item (10 points) Give and justify the recurrence your subproblems should satisfy. 
    \begin{solution}
        \vspace{0.5cm}

        Let's note the $i$-th subproblem as $B[i]$, 
        
        i.e. $B[i]$ is the maximum weight that we can achieve with the first $i$ pumpkins
        
        \vspace{0.4cm}

        We have:

        \[B[1] = A[1]\]
        \[B[2] = \text{max}(A[1],A[2])\]

        Recurrence ralations:

        \[B[i] = \text{max}(B[i-1],(B[i-2]+A[i]))\]

        Pf: 

        \quad For the first \(i\) pumpkins, an optimal selection either 
        
        (1) does not include pumpkin \(i\), giving total \(B[i-1]\); or 
        
        (2) includes pumpkin \(i\), in which case pumpkin \(i-1\) cannot be taken and the total is \(B[i-2]+A[i]\)
        
        The maximum of these two values yields \(B[i]\)
        
        The final answer is \(B[n]\)
        
        The recurrence can be computed in \(O(n)\) time (and \(O(1)\) extra space if we keep only the last two values)

        
        




        
    \end{solution}
    \item (10 points) Using this recurrence, give your algorithm in pseudocode.
    The return value of the algorithm should be the maximum pounds of pumpkin.
    %State and justify the runtime of your algorithm. 
    \begin{solution}
        % \vspace{0.5cm}

        \begin{verbatim}
        Algorithm MaxPumpkinHarvest(A[1..n]):
        
        if n == 0:
            return 0
        if n == 1:
            return A[1]
        
        prev2 = A[1]
        prev1 = max(A[1], A[2])
        
        if n == 2:
            return prev1
        
        for i in range(3,n):
            current = max(prev1, prev2 + A[i])
        prev2 = prev1
        prev1 = current
        
        return prev1
        \end{verbatim}
        

        
    \end{solution}

    
    \item (8 points) State and justify the runtime of your algorithm.
    \begin{solution}
    \vspace{0.5cm}
    
    The runtime of the algorithm is \[ \mathcal{O}(n) \]
    
    \textbf{Justification}
    
    The algorithm processes the array \( A \) of size \( n \) in a single pass
    
    After handling the base cases (which take constant time), it iterates from index 3 to \( n \), 
    
    Performing a constant amount of work in each iteration (a comparison, an addition, and a couple of assignments)
    
    
    Therefore, the total number of operations is proportional to \( n \), yielding a time complexity of \( O(n) \)
    
    
    \end{solution}
    
    % \newpage
    
    \item (6 points) 
    So far, the algorithm that you developed only computes the \emph{pounds} of the pumpkin that can be harvested, but it does not tell you \emph{which pumpkins} to pick.

    Explain how to modify the algorithm so that at the end, we can also produce the list of pumpkins to harvest.
    You may want to store additional information in the DP table as it is filled out.
    
    (\textbf{Hint:} think about what we did in the LCS problem to extract the string after filling out the table.)
\begin{solution}
\vspace{0.5cm}

Similar to the LCS traceback—is to fill the full \( B[1..n] \) table , then backtrack from \( i = n \) down to 1:

Start at \( i = n \)

While \( i \geq 1 \):

  - If \( i = 1 \), 
  
  \hspace{1cm } include pumpkin 1 if \( B[1] > 0 \) (or always, since it’s the only choice)


  \vspace{0.5cm}
  
  - Otherwise, compare \( B[i] \) with \( B[i-1] \):
  
    - If \( B[i] = B[i-1] \), 
    
    \hspace{1cm} then pumpkin \( i \) was not chosen; decrement \( i \) by 1
    
    - If \( B[i] = B[i-2] + A[i] \) (and \( B[i] \ne B[i-1] \)), 
    
    \hspace{1cm} then pumpkin \( i \) was chosen; add \( i \) to the result list and decrement \( i \) by 2

  
This works because the recurrence guarantees that the optimal solution at \( i \) either comes from \( i-1 \) (skip \( i \)) or \( i-2 \) plus \( A[i] \) (take \( i \))

\vspace{1cm}

Modification to the algorithm:

Instead of using constant space, store the entire \( B[1..n] \) array while computing values iteratively

After filling \( B \), perform the backtracking procedure described above to build the list of selected pumpkins

% This increases space usage to \( O(n) \), but allows reconstruction of the optimal set in \( O(n) \) time.

\end{solution}
\end{enumerate}

\section*{Honors Problem -- Optimal Substructure for Rod Cutting (0 Points, *)}

Suppose a profit-maximizing way to cut a rod of length $n$ is given by a sequence of piece lengths
\[
x_1, x_2, \dots, x_k,
\]
where each $x_i$ is a positive integer and $x_1 + x_2 + \cdots + x_k = n$.

\begin{enumerate}
    \item (5 points) Fill in the blank: Then $x_2, \dots, x_k$ is \underline{\hspace{5cm}}.  
    \newline (\textbf{Hint:} Think about the remaining rod of length $n - x_1$.)
    \begin{solution}
        \vspace{1cm}
    \end{solution}
   
\newpage
    \item (20 points) Prove your claim from part (1). (\textbf{Hint:} Try proof by contradiction.)

    \begin{solution}
        \vspace{3cm}
    \end{solution}
   
\end{enumerate}

\section*{Honors Problem -- little-$o$ (0 Points, *)}
We present another asymptotic notation that is used in the analysis of algorithms: For functions $f, g$, we say $f = o(g)$ (``$f$ is little-$o$ of $g$'') if and only if the ratio $\frac{f(n)}{g(n)}$ goes to zero as $n$ tends to infinity. For example, $\log n = o(n)$ but $n\neq o(n)$. 
\begin{enumerate}
    \item Is $\sqrt{n}=o(n)$? Prove that it is or not.
    \begin{solution}
        \vspace{3cm}
    \end{solution}

\newpage
    
    \item Is it possible that $f=o(g)$ and $g=O(f)$? If yes, give an example. Otherwise prove that it is not possible.
    \begin{solution}
        \vspace{4cm}
    \end{solution}
    \item Show that the worst case number of comparisons needed to merge two sorted lists, each of size $n$, is at least $2n-o(n)$.
        \newline
    (Hint: How many ways can we write down $2n$ distinct numbers as two sorted lists of $n$ numbers each? It may also be useful to look up the \href{https://en.wikipedia.org/wiki/Central_binomial_coefficient#Asymptotic_growth}{\underline{\textbf{central binomial coefficient (click here)}}} for its asymptotic growth.)
    \begin{solution}
        \vspace{6cm}
    \end{solution}
\end{enumerate}

